\documentclass[11pt,a4paper]{article}
\input{../../shared/preamble}
\SpeedHeader{Geometry}{2}

\begin{document}

\SpeedTitleBlock{Daily Geometry Drill \#2}{Henry Koduthore}
\SpeedMeta{Circle equations, tangents and line--circle interactions}{completing the square for circle equations, tangent to a circle, perpendicular distance line--circle test, half-chord length}
\noindent\textit{New toolkit today: Circle Equations (centre $(-g,-f)$, $r=\sqrt{g^2+f^2-c}$) $\cdot$ Tangents $\cdot$ Line--Circle Tests $\cdot$ Chord Lengths.}

\vspace{2pt}
\noindent\textit{Attempt each question before checking answers. Time yourself. No calculators.}

\section*{Section A \quad Rapid Recognition \hfill{\normalfont\small($\leq$15 s each)}}
% ══════════════════════════════════════════════════════════════════════════════

\noindent\textit{Complete the square, read the centre and radius, then write or test circle data on sight.}

\begin{enumerate}

\item Write down the centre of the circle $\;x^2 + y^2 - 6x - 8y + 24 = 0$.

\item For which value of $k$ does $\;x^2+y^2-6x+4y+k=0$ represent a circle of radius $5$?

\item Write down the equation of the circle with centre $(2,-1)$ and radius $5$.

\item Write down the radius of the circle $\;x^2 + y^2 - 4x + 6y - 12 = 0$.

\item The point $(x,4)$ lies on the circle $\;x^2 + y^2 = 25$ with $x > 0$. Find $x$.

\item The radius from the origin to the point $(3,4)$ on $\;x^2+y^2=25$ has what gradient?

\item The tangent to $\;x^2 + y^2 = 25$ at the point $(3,4)$ has what gradient?

\item Write down the equation of the circle with centre $(3,-4)$ that passes through the origin.

\item The centre of $\;x^2+y^2-6x-8y+24=0$ is $(3,4)$. Find the perpendicular distance of this centre from the line $\;4x + 3y - 9 = 0$.

\item A chord of the circle $\;x^2+y^2=25$ is at perpendicular distance $3$ from the origin. Find the length of the chord.

\end{enumerate}

% ══════════════════════════════════════════════════════════════════════════════
\section*{Section B \quad Manipulation Drills \hfill{\normalfont\small($\leq$50 s each)}}
% ══════════════════════════════════════════════════════════════════════════════

\noindent\textit{Combine the new circle tools with last drill's line machinery. Complete the square first, then test.}

\begin{enumerate}

\item Which of the following is the equation of the circle with centre $(3,-2)$ and radius $4$?
\begin{itemize}
\item[A)] $x^2+y^2-6x+4y-3=0$
\item[B)] $x^2+y^2-6x+4y+9=0$
\item[C)] $x^2+y^2+6x-4y-3=0$
\item[D)] $x^2+y^2-6x-4y-3=0$
\end{itemize}

\item The circle $\;x^2+y^2+2x-4y-11=0$ has what radius?
\begin{itemize}
\item[A)] $2$
\item[B)] $3$
\item[C)] $4$
\item[D)] $\sqrt{5}$
\end{itemize}

\item The point $(1,2)$ lies on the circle $\;x^2+y^2=5$. The equation of the tangent to the circle at $(1,2)$ is?
\begin{itemize}
\item[A)] $y = -\frac{1}{2}x + \frac{5}{2}$
\item[B)] $y = 2x$
\item[C)] $y = -2x + 5$
\item[D)] $y = \frac{1}{2}x$
\end{itemize}

\item A circle of form $\;x^2+y^2+2gx+2fy+c=0$ passes through the origin. Deduce $c$.
\begin{itemize}
\item[A)] $c = 1$
\item[B)] $c = 0$
\item[C)] $c = g + f$
\item[D)] $c = -1$
\end{itemize}

\item The perpendicular distance from the centre of $\;x^2+y^2=25$ to the line $\;y = 2x + 3$ is?
\begin{itemize}
\item[A)] $\dfrac{3}{\sqrt{5}}$
\item[B)] $\dfrac{3}{5}$
\item[C)] $3$
\item[D)] $2$
\end{itemize}

\item Which of the following is a tangent to the circle $\;x^2+y^2=4$?
\begin{itemize}
\item[A)] $y = x + 2\sqrt{2}$
\item[B)] $y = x + 4$
\item[C)] $y = 2x + 2$
\item[D)] $y = 2x + 4$
\end{itemize}

\item A chord of the circle $\;x^2+y^2=13$ has length $4\sqrt{3}$. What is the perpendicular distance of the chord from the centre?
\begin{itemize}
\item[A)] $3$
\item[B)] $2$
\item[C)] $\sqrt{5}$
\item[D)] $1$
\end{itemize}

\item The line $\;y = mx + 4$ is a tangent to the circle $\;x^2+y^2=4$. Find the two possible values of $m$.
\begin{itemize}
\item[A)] $\pm 1$
\item[B)] $\pm \sqrt{3}$
\item[C)] $\pm 2$
\item[D)] $\pm \sqrt{2}$
\end{itemize}

\item The circle with centre $(6,a)$ and radius $6$ is tangent to the $y$-axis and passes through the point $(0,8)$. Find $a$.
\begin{itemize}
\item[A)] $1$
\item[B)] $3$
\item[C)] $5$
\item[D)] $8$
\end{itemize}

\item Write down the length of the chord cut from the line $\;y = x + 1$ by the circle $\;x^2+y^2=25$.

\end{enumerate}

% ══════════════════════════════════════════════════════════════════════════════
\section*{Section C \quad Substitution \& Structure \hfill{\normalfont\small($\leq$90 s each)}}
% ══════════════════════════════════════════════════════════════════════════════

\noindent\textit{Hard: hidden tangency, non-obvious chord bisection, and the distance-vs-discriminant choice. No brute substitution.}

\begin{enumerate}

\item A tangent to the circle $x^2+y^2=144$ passes through $(20,0)$ and meets the positive $y$-axis at $T$. What is the $y$-coordinate of $T$? \textit{\small(after TMUA 2017 Paper 1 Q6)}
\begin{itemize}
\item[A)] $12$
\item[B)] $15$
\item[C)] $\tfrac{49}{3}$
\item[D)] $20$
\end{itemize}

\item The point $P(8,6)$ is outside the circle $x^2+y^2=36$. Two tangents from $P$ touch the circle at $A$ and $B$. What is the length $AB$ (the common chord of contact)? \textit{\small(adapted from TMUA 2018 Paper 1 Q7)}
\begin{itemize}
\item[A)] $\tfrac{12\sqrt{13}}{13}$
\item[B)] $\tfrac{24\sqrt{13}}{13}$
\item[C)] $\tfrac{48}{5}$
\item[D)] $\tfrac{72}{5}$
\end{itemize}

\item The point $Q$ moves on the circle $x^2+y^2-2x+4y-20=0$, and $P$ is the point $(13,3)$. What is the least possible distance $PQ$?
\begin{itemize}
\item[A)] $13$
\item[B)] $18$
\item[C)] $12$
\item[D)] $8$
\end{itemize}

\item A circle has equation $x^2+y^2-18x-22y+178=0$. A regular hexagon $ABCDEF$ is inscribed so that all its vertices lie on the circle. What is the area of the hexagon? \textit{\small(adapted from TMUA 2017 Paper 1 Q9)}
\begin{itemize}
\item[A)] $18\sqrt3$
\item[B)] $36\sqrt3$
\item[C)] $\tfrac{75\sqrt3}{2}$
\item[D)] $75$
\end{itemize}

\item For which values of $p$ does $x^2-2px+y^2-6y-p^2+8p+9=0$ represent a (real) circle? \textit{\small(after TMUA 2022 Paper 1 Q2)}
\begin{itemize}
\item[A)] $p<0$ or $p>4$
\item[B)] $-1<p<9$
\item[C)] $0<p<4$
\item[D)] $p<-1$ or $p>9$
\end{itemize}

\item The circle $x^2+y^2-4x+2y-11=0$ and the line $x-y+1=0$: determine the number of intersection points without solving the system. \textit{\small(after MAT 2019 Q6)}
\begin{itemize}
\item[A)] $0$ (line misses)
\item[B)] $1$ (tangent)
\item[C)] $2$ (secant)
\item[D)] It cannot be decided without solving the system
\end{itemize}

\item The circle $x^2+y^2-8x-4y-21=0$ is cut by the $y$-axis. What is the length of the chord intercepted on the $y$-axis?
\begin{itemize}
\item[A)] $2\sqrt{21}$
\item[B)] $2\sqrt{41}$
\item[C)] $8$
\item[D)] $10$
\end{itemize}

\item The line $y=kx+5$ is tangent to the circle $x^2+y^2-6x-8y+24=0$. How many distinct values of $k$ satisfy this? \textit{\small(after TMUA 2020 Paper 1 Q7)}
\begin{itemize}
\item[A)] $0$
\item[B)] $1$
\item[C)] $2$
\item[D)] Infinitely many
\end{itemize}

\end{enumerate}

% ══════════════════════════════════════════════════════════════════════════════
\section*{Section D \quad Challenge \hfill{\normalfont\small(TMUA / SMC / BMO1 difficulty)}}
% ══════════════════════════════════════════════════════════════════════════════

\noindent\textit{Hardest 5: distance invariants, hidden right angles, and a construction you must spot.}

\begin{enumerate}

\item The circles $C_1:(x+2)^2+(y-3)^2=18$ and $C_2:(x-7)^2+(y+6)^2=2$ have centres $O_1,O_2$. What is the \emph{shortest} distance between a point on $C_1$ and a point on $C_2$? \textit{\small(adapted from TMUA 2018 Paper 1 Q7 / 2022 Paper 1 Q3)}
\begin{itemize}
\item[A)] $5\sqrt2-4$
\item[B)] $5\sqrt2-5$
\item[C)] $5\sqrt2$
\item[D)] $5\sqrt2+5$
\end{itemize}

\item The segment joining $(3,3)$ and $(7,5)$ is a diameter of circle $C$. $C$ is translated $3$ left, reflected in the $x$-axis, then enlarged scale factor $4$ about its centre. What is the equation of the final circle? \textit{\small(after TMUA Specimen Paper 1 Q9)}
\begin{itemize}
\item[A)] $(x-2)^2+(y-4)^2=80$
\item[B)] $(x-2)^2+(y+4)^2=80$
\item[C)] $(x-2)^2+(y-4)^2=320$
\item[D)] $(x-2)^2+(y+4)^2=320$
\end{itemize}

\item A circle has centre $O$ and radius $6$. Points $P$ and $Q$ lie on it with $\angle POQ\ge 90^\circ$, and triangle $POQ$ has area $9\sqrt3$. The point $R$ moves on the major arc $PQ$. What is the greatest possible area of triangle $PRQ$? \textit{\small(adapted from TMUA 2022 Paper 1 Q14)}
\begin{itemize}
\item[A)] $18+9\sqrt3$
\item[B)] $27\sqrt3$
\item[C)] $27+9\sqrt3$
\item[D)] $36+9\sqrt3$
\end{itemize}

\item Two tangents are drawn from the origin to the circle $(x-5)^2+(y-5)^2=5$. What is the tangent of the acute angle between them?
\begin{itemize}
\item[A)] $\tfrac34$
\item[B)] $\tfrac43$
\item[C)] $\tfrac52$
\item[D)] $1$
\end{itemize}

\item The circle through $(0,0)$, $(8,0)$ and $(2,6)$ meets the $y$-axis again at $Q$. How far is $Q$ from the origin?
\begin{itemize}
\item[A)] $2\sqrt5$
\item[B)] $6$
\item[C)] $2$
\item[D)] $4$
\end{itemize}

\end{enumerate}

\SpeedClosing{``Tangency is a distance statement in disguise: compare $d$ to $r$, never solve the intersection.''}

\end{document}
